\chapter*{\Large \center Abstract}

Visitors to Sheffield's Festival of the Mind typed how they felt and a text-to-image installation answered each entry with a monochrome line drawing. That left \nTotal{} unlabelled prompt--drawing records, of which I analyse \nAnalysed{}. Nobody rated the drawings, so there is no ground truth. I ask instead whether one model's reading of a drawing agrees with another model's reading of its prompt. The prompt reader is RoBERTa fine-tuned on GoEmotions. The image reader is describe-then-extract, where BLIP-VQA is asked five factual questions and BART-MNLI reads an emotion family from the resulting description. On my supervisor's advice the readings are scored on a four-rung similarity ladder, from valence--arousal similarity to family agreement, and never matched label to label. Because valence--arousal similarity cannot fall below 0.412, with random pairs averaging \FullNeuVaNull{}, every score is reported beside its random-pairing null, and because \nRepeats{} records repeat a prompt, the intervals and tests are clustered on prompt groups.

On the \nEmotional{} prompts the prompt reader labels emotional, observed valence-sign agreement is \FullEmoSignObs{} where random pairing averages \FullEmoSignNull{} ($z$ = \FullEmoSignZ{}), and family agreement is \FullEmoFamObs{} against \FullEmoFamNull{}. Both hold with each prompt counted once and after Holm correction. Under the describer's split, face-classified drawings reach \FullFaceSignObs{} on valence sign and scene-classified drawings reach \FullSceneSignObs{}, a gap of \DiffSignPp{} (clustered 95\% interval \DiffSignCi{}, clustered permutation $p$ = \DiffSignP{}). Neutral-labelled prompts stay at chance. I also ran direct CLIP like for like; it gets \McClipFam{} family agreement where describe-then-extract gets \McDeFam{}. Where a prompt repeats, \ConsFam{} of its drawings receive the group's modal family, against \ConsFamNull{} with families shuffled.

The describer reports no person in \DescScenePeopleZero{} scene-classified drawings yet gives an expression for \DescScenePeopleZeroExpr{} of them, and corpus-wide, wherever that answer is an emotion word, the image reader returns that word's family \DescCopyStrict{} of the time. Dropping the expression sentence cuts valence-sign agreement to \AblDropExpressionEmoSignObs{} (null \AblDropExpressionEmoSignNull{}), and a people-aware template leaves the face rows unchanged and the scene rows at chance, so the headline agreement is a property of the measurement chain. To say what a person sees, I would need the single-coder check, then a human-rated study. The check was unscored when I wrote this. The dissertation contributes the corpus and the research question, the null-referenced ladder, the consistency experiment, and an audited, reproducible toolkit and demonstrator.
